\begin{figure}[H]
\centering
\begin{tikzpicture}[font=\small,>=Stealth]

\draw[gray] (0,0) circle (1.25);
\draw[blue,very thick] (1.25,0) arc (0:180:1.25);
\draw[teal,very thick] (0,-1.25) arc (-90:90:1.25);
\node at (0,1.75) {$U$: upper semicircle};
\node[teal] at (1,-1.65) {$V$: right semicircle};
\draw[blue,thick] (4,1)--(8,1);
\draw[teal,thick] (4,-1)--(8,-1);
\node[right] at (8,1) {$x\in(-1,1)$};\node[right] at (8,-1) {$y\in(-1,1)$};
\draw[->,blue] (1,1)--(3.8,1) node[midway,above] {$\varphi$};
\draw[->,teal] (1.2,-.3)--(3.8,-1) node[midway,above] {$\psi$};
\draw[->,thick] (6,.8)--(6,-.8) node[midway,right] {$y=\sqrt{1-x^2}$};
\node[align=center] at (4,-2.3) {on the upper-right overlap, $0<x,y<1$};

\end{tikzpicture}
\caption{Two ways to flatten circle neighborhoods into intervals. The transition compares coordinates of the same point on their overlap; the inverse chart returns a coordinate to the circle.}
\label{fig:preview-charts}
\end{figure}
